\chapter{Vecteurs, translations et combinaisons}
\section*{Exercice 8.1 -- Déplacements}
$A(2,-1)$, $B(5,3)$, $C(-4,2)$. Calculer $\overrightarrow{AB}$ puis le point $D$ tel que $\overrightarrow{CD}=\overrightarrow{AB}$.

\section*{Exercice 8.2 -- Combinaison}
Pour $\vec u=(2,-3)$ et $\vec v=(-1,4)$, calculer $3\vec u-2\vec v$ et sa norme dans une base orthonormée.

\section*{Exercice 8.3 -- Chasles}
Simplifier $\overrightarrow{AB}+\overrightarrow{BC}-\overrightarrow{DC}$.

\section*{Exercice 8.4 -- Parallélogramme}
$A(1,1)$, $B(4,2)$, $C(5,6)$ sont trois sommets successifs d'un parallélogramme $ABCD$. Trouver $D$.

\section*{Exercice 8.5 -- Milieu vectoriel}
Montrer que $\overrightarrow{MA}+\overrightarrow{MB}=2\overrightarrow{MI}$ lorsque $I$ est le milieu de $[AB]$.

\section*{Exercice 8.6 -- Changement d'échelle}
Pour un vecteur non nul $\vec u$ de norme $5$, décrire $-\tfrac32\vec u$ : direction, sens et norme. Prouver la formule de norme par coordonnées.
